\begin{ex}
Gấp 6 lên 5 lần ta được số nào?
\shortans{30}
\loigiai{
$6 \times 5 = 30$.
}
\end{ex}

\begin{ex}
Gấp 8 lên 4 lần ta được số nào?
\shortans{32}
\loigiai{
$8 \times 4 = 32$.
}
\end{ex}

\begin{ex}
Gấp 9 lên 7 lần ta được số nào?
\shortans{63}
\loigiai{
$9 \times 7 = 63$.
}
\end{ex}

\begin{ex}
Gấp 12 lên 3 lần ta được số nào?
\shortans{36}
\loigiai{
$12 \times 3 = 36$.
}
\end{ex}

\begin{ex}
Gấp 15 lên 4 lần ta được số nào?
\shortans{60}
\loigiai{
$15 \times 4 = 60$.
}
\end{ex}

\begin{ex}
Gấp 16 lên 5 lần ta được số nào?
\shortans{80}
\loigiai{
$16 \times 5 = 80$.
}
\end{ex}

\begin{ex}
Gấp 24 lên 3 lần ta được số nào?
\shortans{72}
\loigiai{
$24 \times 3 = 72$.
}
\end{ex}

\begin{ex}
Gấp 18 lên 4 lần ta được số nào?
\shortans{72}
\loigiai{
$18 \times 4 = 72$.
}
\end{ex}

\begin{ex}
Đoạn thẳng AB dài 7 cm. Đoạn thẳng CD dài gấp 4 lần đoạn thẳng AB. Đoạn thẳng CD dài bao nhiêu xăng-ti-mét?
\shortans{28}
\loigiai{
$7 \times 4 = 28$ cm.
}
\end{ex}

\begin{ex}
Bạn Lan có 6 cái nhãn vở. Bạn Mai có số nhãn vở gấp 3 lần số nhãn vở của bạn Lan. Hỏi bạn Mai có bao nhiêu cái nhãn vở?
\shortans{18}
\loigiai{
$6 \times 3 = 18$ cái.
}
\end{ex}

\begin{ex}
Năm nay em 6 tuổi. Tuổi mẹ gấp 5 lần tuổi em. Hỏi năm nay mẹ bao nhiêu tuổi?
\shortans{30}
\loigiai{
$6 \times 5 = 30$ tuổi.
}
\end{ex}

\begin{ex}
Năm nay con 7 tuổi, tuổi bố gấp 5 lần tuổi con. Hỏi bố bao nhiêu tuổi?
\shortans{35}
\loigiai{
$7 \times 5 = 35$ tuổi.
}
\end{ex}

\begin{ex}
Bao gạo thứ nhất nặng 12 kg. Bao gạo thứ hai nặng gấp 3 lần bao thứ nhất. Hỏi bao thứ hai nặng bao nhiêu ki-lô-gam?
\shortans{36}
\loigiai{
$12 \times 3 = 36$ kg.
}
\end{ex}

\begin{ex}
Một con vịt nặng 2 kg, một con lợn nặng gấp 25 lần con vịt. Hỏi con lợn nặng bao nhiêu ki-lô-gam?
\shortans{50}
\loigiai{
$2 \times 25 = 50$ kg.
}
\end{ex}

\begin{ex}
Hùng có 8 viên bi, số bi của Dũng gấp 6 lần số bi của Hùng. Hỏi Dũng có bao nhiêu viên bi?
\shortans{48}
\loigiai{
$8 \times 6 = 48$ viên bi.
}
\end{ex}

\begin{ex}
Một số gấp lên 6 lần thì được 42. Số đó là số nào?
\shortans{7}
\loigiai{
$42 : 6 = 7$.
}
\end{ex}

\begin{ex}
Một số gấp lên 8 lần thì được 72. Số đó là số nào?
\shortans{9}
\loigiai{
$72 : 8 = 9$.
}
\end{ex}

\begin{ex}
Một số gấp lên 4 lần rồi cộng với 12 thì được 48. Số đó là số nào?
\shortans{9}
\loigiai{
Số đó gấp 4 lần: $48 - 12 = 36$. Số đó là $36 : 4 = 9$.
}
\end{ex}

\begin{ex}
Thùng thứ nhất có 15 lít dầu. Thùng thứ hai có số lít dầu gấp 4 lần thùng thứ nhất. Hỏi thùng thứ hai có bao nhiêu lít dầu?
\shortans{60}
\loigiai{
$15 \times 4 = 60$ lít.
}
\end{ex}

\begin{ex}
Mẹ hái được 8 quả bưởi, bà hái được gấp 3 lần số bưởi của mẹ. Hỏi cả hai người hái được bao nhiêu quả bưởi?
\shortans{32}
\loigiai{
Bà hái: $8 \times 3 = 24$ quả. Cả hai người: $8 + 24 = 32$ quả.
}
\end{ex}

\begin{ex}
Gấp 14 lên 5 lần rồi bớt đi 20 thì được bao nhiêu?
\shortans{50}
\loigiai{
$14 \times 5 - 20 = 70 - 20 = 50$.
}
\end{ex}

\begin{ex}
Gấp đôi của 35 là bao nhiêu?
\shortans{70}
\loigiai{
$35 \times 2 = 70$.
}
\end{ex}

\begin{ex}
Gấp ba của 24 là bao nhiêu?
\shortans{72}
\loigiai{
$24 \times 3 = 72$.
}
\end{ex}

\begin{ex}
Gấp bốn của 16 là bao nhiêu?
\shortans{64}
\loigiai{
$16 \times 4 = 64$.
}
\end{ex}

\begin{ex}
Một hình vuông có cạnh 9 cm. Nếu gấp cạnh của nó lên 2 lần thì cạnh mới dài bao nhiêu xăng-ti-mét?
\shortans{18}
\loigiai{
$9 \times 2 = 18$ cm.
}
\end{ex}
